% LSPV-2 REPORT::2021::2026-09-30T20:21:35-07:00
% !TEX program = xelatex
\documentclass[11pt]{article}
\usepackage[margin=0.9in]{geometry}
\usepackage{amsmath,amssymb}
\begin{document}
\section*{LSPV-2: accepted-family cardinality obstruction to low product span}
\textbf{STATUS: OPEN.} P versus NP remains open. This run does not prove a rank bound for the final SAME\_REALIZATION family.
\paragraph{Sources.} Reviewed repository protocol/skill, root \texttt{main.tex}, MAIN 1958, LSPV-1 LEAF2/SCHUR1, LSPV-2 1723, LSPV-3 1833A-R1, LSPV-4 1844R, GROK BOOL v2.8, and CLAUDE-0266. Primary literature checked: Shaltiel--Umans, \emph{SIAM J. Comput.} 39 (2009), and van Melkebeek--Mocelin Sdroievski, CCC 2023.
\paragraph{Target and quantifiers.} Fix one successful vMS SAME\_REALIZATION path and one RM node at which a fixed local extractor \(E\) is applied. Let \(\mathcal F_{\rm acc}\) be the set of distinct reachable oracle objects at that node induced by perfectly accepted PCP proofs on this same path. For \(f\in\mathcal F_{\rm acc}\), let
\[
H_f=\{E^f(y):y\in\{0,1\}^d\}\subseteq\mathbb F_2^m,\qquad M=1+m+\binom m2,
\]
and let \(\rho(f)\) be the degree-\(\le2\) feature rank of \(H_f\).
\paragraph{PROVED/EXACT: accepted-family cardinality obstruction.} LSPV-1 LEAF2 derived from Shaltiel--Umans Proposition 3.8 that, for a \((\kappa,\varepsilon)\) local extractor with \(\varepsilon<1/4\),
\[
\left|\{f\in\mathcal C:\rho(f)<M\}\right|<2^{\kappa+M}.
\]
Intersecting with the reachable family gives
\[
\boxed{|\mathcal F_{\rm acc}|>2^{\kappa+M}\Longrightarrow \exists f\in\mathcal F_{\rm acc}\text{ with }\rho(f)=M.}
\]
Since the final SAME\_REALIZATION range contains \(H_f\), its quadratic-feature rank is then \(M\), so \(K_2(H_{\rm final})=\{0\}\). Thus the quadratic sparse-\(Q\) route is retired for that realization whenever this threshold is met. Equivalently, a necessary condition for global rank deficiency is
\[
\boxed{|\mathcal F_{\rm acc}|\le 2^{\kappa+M}}
\]
at every extractor node whose output is included in the final range.
\paragraph{SAME\_REALIZATION significance.} The broad interpolation counterexample did not prove reachability from perfectly accepted PCP proofs. The new criterion asks only for enough distinct reachable accepted objects. Many reachable accepted objects force full quadratic rank, while a low-rank proof must establish a strong upper bound after the exact LDE/restriction map.
\paragraph{Important non-implications.} Proof multiplicity alone is insufficient. LSPV-4's padding control can create many accepting proof strings with identical verifier semantics. If they collapse to the same node oracle object, \(|\mathcal F_{\rm acc}|\) does not increase. The relevant cardinality is after \(y\mapsto\widehat g_y\mapsto f_{\rm node}\), with equality tested extensionally at the oracle interface. This blocks decorative proof padding, COPY wrappers, and source relabeling.
At univariate leaves the criterion is presently vacuous: LSPV-1 computed \(\kappa=\Delta^5\), while the entire degree-\(\le\Delta\) univariate class has only \(p^{\Delta+1}\) objects and \(\log_2p^{\Delta+1}=O(\Delta\log\Delta)\ll\Delta^5\). The criterion must be tested at higher-dimensional RM nodes.
\paragraph{Sparse-\(Q\) accounting.} If instead a construction-specific argument proves \(\rho=r_2<M\), Singleton yields a nonzero quadratic \(Q\) with at most \(r_2+1\) ANF monomials, \(O(r_2\log m)\) monomial-list bits, and \(O(r_2)\) direct Boolean evaluation gates. Acquisition, exact size normalization, census/output, and reconstruction remain separate charges.
\paragraph{Hostile/barrier audit.} The obstruction is invariant under proof-string relabeling but not wrappers altering the extensional node oracle; such changes must be SAME\_REALIZATION legal. Horn, 2-SAT, affine/Tseitin, and bounded-treewidth controls forbid interpreting high rank as runtime hardness. A maximal semantic or algebraic oracle can expose rank and accepted-object identities, so this argument relativizes and is not the missing hinge. It supplies no arbitrary-algorithm transfer.
\paragraph{Handoffs.} LSPV-1: at the first higher-dimensional RM node, compute the exact map from perfectly accepted PCP proofs to distinct oracle objects and compare \(\log_2|\mathcal F_{\rm acc}|\) with \(\kappa+M\). LSPV-3: no reconstruction improvement follows unless rank deficiency survives. LSPV-4: attack via collisions in the accepted-proof-to-node-oracle map. LSPV-5: the quadratic route now has a sharp reachable-family falsification threshold.
\paragraph{Single next target.} For the earliest higher-dimensional RM node on one successful Theorem 24 path, determine whether \(\log_2|\mathcal F_{\rm acc}|>\kappa+1+m+\binom m2\). A proved strict inequality gives an architecture-preserving full-rank witness and retires quadratic contraction for that realization; failure requires identifying the exact collision/compression mechanism.
\begin{thebibliography}{9}
\bibitem{SU09} R. Shaltiel and C. Umans, Low-End Uniform Hardness versus Randomness Tradeoffs for AM, \emph{SIAM J. Comput.} 39(3):1006--1037, 2009, doi:10.1137/070698348.
\bibitem{vMS23} D. van Melkebeek and N. Mocelin Sdroievski, Instance-Wise Hardness Versus Randomness Tradeoffs for Arthur-Merlin Protocols, CCC 2023, LIPIcs 264:17, doi:10.4230/LIPIcs.CCC.2023.17.
\end{thebibliography}
\end{document}
% END_OF_VERIFIED_REPORT::LSPV-2::2021